\chapter{Introduction}
In recent years, an interesting challenge in the field of Artificial Intelligence (AI) and safety-critical systems has been the integration of AI models into systems that require predictable and safe behaviors. The black-box nature of AI models implies an absence of structured standards and methodologies to obtain AI-based safety-critical systems.

In this project, I will begin the study of these topics by examining one of the key papers~\cite{perez2024artificial} in this discipline.

\myskip

This is the \href{https://github.com/edoardosarri24/AI-Safety-Critical-runtime-monitoring}{repository} of the project.

%%%%%%%%%%%%%%%%%%%%%%%%%%%%%%%%%%%%%%%%%%%%%%%%
\section{Tools used}
During the development of this project, in order to improve the report, the development process, and the quality of the overall work, I used the following tools:
\begin{itemize}
    \item \textbf{NotebookLM} \\
        It was useful for performing an initial review of the guideline paper~\cite{perez2024artificial}. \href{https://notebooklm.google.com/notebook/26530d12-ce92-4268-a231-5d9a3e5e7358}{Here} is the link to the notebook: although the chat history is no longer available, you can see the main interactions, such as the concept map and report.
    \item \textbf{Copilot completions} \\
        While writing this report (and the slides), Copilot completions generated some parts of the text for me.
    \item \textbf{Gemini CLI} \\
        The CLI version of Gemini was used to check the grammatical correctness of my English after writing.
\end{itemize}

%%%%%%%%%%%%%%%%%%%%%%%%%%%%%%%%%%%%%%%%%%%%%%%%
\section{Terms}
In this short section, we clarify some of the terms used in the remainder of this work.
\begin{itemize}
    \item \textbf{AI and Machine Learning} \\
        This refers to the algorithmic techniques that allow us to create models that learn patterns from data.
    \item \textbf{AI Safety} \\
        This is used to refer to the techniques used to avoid or mitigate the potential catastrophic harm of AI systems to humans. This means that AI safety also refers to the methodologies that allow us to deploy AI systems in safety-critical systems while being compliant with safety standards.
    \item \textbf{Functional Safety (FuSa)} \\
        Functional Safety is defined as the "part of the overall safety of a system that assures the freedom from unacceptable risk" [120]~\cite{perez2024artificial}. The aim of FuSa is to ensure that a system is free from unacceptable risk.
\end{itemize}